\section{Agenda de Pesquisa}

A RQ5 pergunta quais lacunas permanecem para comparar, avaliar e governar esses frameworks em projetos reais. Diferente das demais perguntas, ela é respondida no nível do corpus inteiro, e não estudo a estudo. Da síntese emergem seis lacunas, cada uma com uma direção de pesquisa concreta.

\noindent\textbf{Ausência de benchmarks orientados a processo.} A avaliação do corpus concentra-se no resultado final, isto é, no código gerado e na taxa de resolução de issues, e não na qualidade do processo: aderência da implementação à especificação, taxa de desalinhamento, custo de coordenação entre agentes e qualidade dos artefatos intermediários. Nenhum estudo propõe um benchmark padronizado de processo. A direção indicada é definir métricas e conjuntos de dados que avaliem o pipeline, e não apenas a solução.

\noindent\textbf{Falta de comparação empírica direta entre frameworks.} Cada estudo avalia o próprio framework isoladamente. Não há estudo comparativo controlado que confronte, por exemplo, ChatDev, MetaGPT e AgileGen na mesma tarefa, com as mesmas métricas \cite{ChenQian2023,SiruiHong2023,SaiZhang2025}. Isso bloqueia conclusões sobre qual estrutura de processo funciona melhor e sob quais condições. A direção indicada são experimentos comparativos diretos com protocolos comuns.

\noindent\textbf{Portabilidade e aprisionamento, pouco estudados como objeto.} A dependência de IDE, plataforma e modelo comercial é quase universal nas fichas, mas apenas um estudo nomeia o aprisionamento explicitamente, e nenhum mede portabilidade de forma sistemática \cite{SeyedmoeinMohsenimofidi2025}. A direção indicada são instrumentos para medir o acoplamento a fornecedor e o custo de migração entre plataformas agênticas.

\noindent\textbf{Governança e rastreabilidade carecem de validação em escala real.} As propostas de traço, contrato e auditoria são promissoras, mas validadas em estudos de caso ou protótipos, não em organizações ao longo do tempo, e falta análise de custo-benefício do custo de coordenação imposto pela governança \cite{Paduraru2026,Merchant2025,Watfa2026,Macedo2026}. A direção indicada é a avaliação longitudinal de regimes de governança em equipes reais.

\noindent\textbf{A dimensão humana e de adoção está sub-representada na taxonomia técnica.} O tema emergente de colaboração mostra que confiança e adoção organizacional são decisivos, mas a taxonomia de seis dimensões não os contempla \cite{DanielRusso2024,Ulfsnes2024,Baron2025,Akbar2025}. A direção indicada é integrar uma camada sociotécnica à classificação de frameworks e estudar fatores de adoção em campo.

\noindent\textbf{Segurança tratada reativamente.} Com a exceção de ASTRA, segurança aparece como risco a mitigar, e não como propriedade de projeto avaliada \cite{XiangzheXu2025}. Falta avaliação sistemática da superfície de ataque de frameworks agênticos, em especial de extensões, skills e plugins comunitários. A direção indicada é um framework de avaliação de segurança específico para pipelines de desenvolvimento agêntico.
